\section{Restricciones del sistema y garantías de estabilidad}
\subsection{Límites de recursos y de cómputo}
El entrenamiento de gradiente de política suele enfrentar un cómputo y un budget limitados: las trayectorias de horizon largo implican un tiempo de GPU enorme, mientras que la entropy regularization requiere compute adicional. Solemos equilibrar la carga de entrenamiento mediante accumulate gradients, recorte de gradiente y mixed precision, y agregamos verificaciones de dependency temporal en el pipeline para evitar drift.

\begin{table}[H]
\centering
\small
\begin{tabular}{p{0.34\textwidth}p{0.38\textwidth}p{0.28\textwidth}}
\toprule
Tipo de restricción & Manifestación concreta & Respuesta de ingeniería \\
\midrule
Latencia/presupuesto & La horizon larga de la trayectoria provoca un compute burst & gradient accumulation + distributed rollout \\
Estabilidad & El entropy collapse provoca policy collapse & entropy bonus + KL guard \\
Disponibilidad de datos & limited rollout data per day & replay buffer + prioritized sampling \\
\bottomrule
\end{tabular}
\caption{Panorama de recursos del gradiente de política bajo restricciones del sistema}
\end{table}

\begin{importantbox}{Vínculo entre el cómputo y el logging}
\begin{itemize}
    \item En cada etapa de rollout deben registrarse seed, config y entropy curve para permitir la reproducción;
    \item Las trayectorias de horizon largo requieren chunking al registrarse, para evitar que el archivo de log se dispare;
    \item En el closed-loop pipeline se vincula el gradient step count con el runtime metric, lo que facilita el capacity planning.
\end{itemize}
\end{importantbox}

\subsection{Institucionalidad y equidad}
Además del hardware, el despliegue del modelo también está sujeto a restricciones institucionales de cumplimiento, equidad y seguridad. Por ejemplo, en un multi-domain rollout es necesario detectar la distribución del reward de la policy entre distintos grupos; si la brecha es demasiado grande, es necesario trigger guardrail. Aquí se suele usar un fairness dashboard para registrar disparity, reject rate y manual override logs.

\begin{warningbox}{Riesgos de ignorar los requisitos institucionales}
Si solo se presta atención al reward y se ignoran las métricas de fairness, la policy podría ser detenida por el regulator tras el release debido a un desempeño diferenciado, lo que retrasaría todo el avance del producto.
\end{warningbox}

\begin{knowledgebox}{Prácticas de garantía institucionalizada}
\begin{itemize}
    \item Establecer un `fairness gate`: realizar una evaluación anticipada del reward/distribution de los subgroup sensibles;
    \item Aplicar structured logging a los `override event`, vinculándolos con la policy version y el prompt template;
    \item Usar un audit-ready benchmark (como fairness-specific tasks) para garantizar la alineación en múltiples dimensiones.
\end{itemize}
\end{knowledgebox}

\subsection{Resumen de la sección}
Las restricciones del sistema provienen tanto del budget de cómputo como de los requisitos de equidad y compliance. Al refinar el logging, el entropy guard y el fairness gate, es posible llevar el entrenamiento de gradiente de política de regreso a una vía de ingeniería manejable.

